\section{Cohen-Macaulay modules}
\label{section-CM}

\noindent
Here we show that Cohen-Macaulay modules have good properties. We postpone
using Ext groups to establish the connection with duality and so on.

\begin{definition}
\label{definition-CM}
Let $R$ be a Noetherian local ring.
Let $M$ be a finite $R$-module.
We say $M$ is {\it Cohen-Macaulay}
if $\dim(\text{Supp}(M)) = \text{depth}(M)$.
\end{definition}

\noindent
A first goal will be to establish Proposition \ref{proposition-CM-module}.
We do this by a (perhaps nonstandard) sequence of elementary lemmas
involving almost none of the earlier results on depth. Let us introduce
some notation.

\medskip\noindent
Let $R$ be a local Noetherian ring. Let $M$ be
a Cohen-Macaulay module, and let $f_1, \ldots, f_d$
be an $M$-regular sequence with $d = \dim(\text{Supp}(M))$.
We say that $g \in \mathfrak m$ is {\it good with respect to
$(M, f_1, \ldots, f_d)$} if for all $i = 0, 1, \ldots, d-1$
we have $\dim (\text{Supp}(M) \cap V(g, f_1, \ldots, f_i))
= d - i - 1$. This is equivalent to the condition that
$\dim(\text{Supp}(M/(f_1, \ldots, f_i)M) \cap V(g)) =
d - i - 1$ for $i = 0, 1, \ldots, d - 1$.

\begin{lemma}
\label{lemma-good-element}
Notation and assumptions as above. If $g$ is good with respect to
$(M, f_1, \ldots, f_d)$, then (a) $g$ is a nonzerodivisor on $M$,
and (b) $M/gM$ is Cohen-Macaulay with maximal regular
sequence $f_1, \ldots, f_{d - 1}$.
\end{lemma}

\begin{proof}
We prove the lemma by induction on $d$.
If $d = 0$, then $M$ is finite and there is no case
to which the lemma applies.
If $d = 1$, then we have to show that $g : M \to M$ is
injective. The kernel $K$ has support $\{\mathfrak m\}$
because by assumption $\dim \text{Supp}(M) \cap V(g) = 0$.
Hence $K$ has finite length. Hence $f_1 : K \to K$ injective
implies the length of the image is the length of $K$, and hence
$f_1 K = K$, which by Nakayama's Lemma \ref{lemma-NAK} implies $K = 0$.
Also, $\dim \text{Supp}(M/gM) = 0$ and so $M/gM$ is Cohen-Macaulay
of depth $0$.

\medskip\noindent
Assume $d > 1$. Observe that $g$ is good for $(M/f_1M, f_2, \ldots, f_d)$,
as is easily seen from the definition. By induction, we have that
(a) $g$ is a nonzerodivisor on $M/f_1M$ and
(b) $M/(g, f_1)M$ is Cohen-Macaulay with maximal regular sequence
$f_2, \ldots, f_{d - 1}$. By
Lemma \ref{lemma-permute-xi}
we see that $g, f_1$ is an $M$-regular sequence.
Hence $g$ is a nonzerodivisor on $M$ and
$f_1, \ldots, f_{d - 1}$ is an $M/gM$-regular sequence.
\end{proof}

\begin{lemma}
\label{lemma-CM-one-g}
Let $R$ be a Noetherian local ring.
Let $M$ be a Cohen-Macaulay module over $R$.
Suppose $g \in \mathfrak m$ is such that $\dim(\text{Supp}(M) \cap V(g))
= \dim(\text{Supp}(M)) - 1$. Then (a) $g$ is a nonzerodivisor on $M$,
and (b) $M/gM$ is Cohen-Macaulay of depth one less.
\end{lemma}

\begin{proof}
Choose a $M$-regular sequence $f_1, \ldots, f_d$ with
$d = \dim(\text{Supp}(M))$. If $g$ is good with respect to
$(M, f_1, \ldots, f_d)$ we win by Lemma \ref{lemma-good-element}.
In particular the lemma holds if $d = 1$. (The case $d = 0$ does
not occur.) Assume $d > 1$. Choose an element $h \in R$ such that
(\romannumeral1) $h$ is good with respect to $(M, f_1, \ldots, f_d)$,
and (\romannumeral2) $\dim(\text{Supp}(M) \cap V(h, g)) = d - 2$.
To see $h$ exists, let $\{\mathfrak q_j\}$ be the (finite) set of
minimal primes of the closed sets $\text{Supp}(M)$,
$\text{Supp}(M)\cap V(f_1, \ldots, f_i)$, $i = 1, \ldots, d - 1$,
and $\text{Supp}(M) \cap V(g)$. None of these $\mathfrak q_j$
is equal to $\mathfrak m$ and hence we may find $h \in \mathfrak m$,
$h \not \in \mathfrak q_j$ by Lemma \ref{lemma-silly}. It is clear
that $h$ satisfies (\romannumeral1) and (\romannumeral2). From
Lemma \ref{lemma-good-element} we conclude that
$M/hM$ is Cohen-Macaulay. By (\romannumeral2) we see that the pair
$(M/hM, g)$ satisfies the induction hypothesis. Hence
$M/(h, g)M$ is Cohen-Macaulay and $g : M/hM \to M/hM$
is injective. By Lemma \ref{lemma-permute-xi} we see
that $g : M \to M$ and $h : M/gM \to M/gM$
are injective. Combined with the fact that $M/(g, h)M$
is Cohen-Macaulay this finishes the proof.
\end{proof}

\begin{proposition}
\label{proposition-CM-module}
Let $R$ be a Noetherian local ring, with maximal ideal $\mathfrak m$.
Let $M$ be a Cohen-Macaulay module over $R$ whose support has dimension $d$.
Suppose that $g_1, \ldots, g_c$ are elements of
$\mathfrak m$ such that $\dim(\text{Supp}(M/(g_1, \ldots, g_c)M))
= d - c$. Then $g_1, \ldots, g_c$ is an $M$-regular sequence,
and can be extended to a maximal $M$-regular sequence.
\end{proposition}

\begin{proof}
Let $Z = \text{Supp}(M) \subset \Spec(R)$.
By Lemma \ref{lemma-one-equation} in the chain
$Z \supset Z \cap V(g_1) \supset \ldots \supset Z \cap V(g_1, \ldots, g_c)$
each step decreases the dimension at most by $1$. Hence by assumption
each step decreases the dimension by exactly $1$ each time. Thus we
may successively apply Lemma \ref{lemma-CM-one-g} to the modules
$M/(g_1, \ldots, g_i)$ and the element $g_{i + 1}$.

\medskip\noindent
To extend $g_1, \ldots, g_c$ by one element if $c < d$ we simply
choose an element $g_{c + 1} \in \mathfrak m$ which is not
in any of the finitely many minimal primes of $Z \cap V(g_1, \ldots, g_c)$,
using Lemma \ref{lemma-silly}.
\end{proof}

\noindent
Having proved Proposition \ref{proposition-CM-module} we continue the
development of standard theory.

\begin{lemma}
\label{lemma-nonzerodivisor-on-CM}
Let $R$ be a Noetherian local ring with maximal ideal $\mathfrak m$.
Let $M$ be a finite $R$-module. Let $x \in \mathfrak m$ be a
nonzerodivisor on $M$. Then $M$ is Cohen-Macaulay if and only
if $M/xM$ is Cohen-Macaulay.
\end{lemma}

\begin{proof}
By Lemma \ref{lemma-depth-drops-by-one} we have
$\text{depth}(M/xM) = \text{depth}(M)-1$.
By Lemma \ref{lemma-one-equation-module}
we have $\dim(\text{Supp}(M/xM)) = \dim(\text{Supp}(M)) - 1$.
\end{proof}

\begin{lemma}
\label{lemma-CM-over-quotient}
Let $R \to S$ be a surjective homomorphism of Noetherian local rings.
Let $N$ be a finite $S$-module. Then $N$ is Cohen-Macaulay as an $S$-module
if and only if $N$ is Cohen-Macaulay as an $R$-module.
\end{lemma}

\begin{proof}
Omitted.
\end{proof}

\begin{lemma}
\label{lemma-CM-ass-minimal-support}
\begin{reference}
\cite[Chapter 0, Proposition 16.5.4]{EGA}
\end{reference}
Let $R$ be a Noetherian local ring. Let $M$ be a finite Cohen-Macaulay
$R$-module. If $\mathfrak p \in \text{Ass}(M)$, then
$\dim(R/\mathfrak p) = \dim(\text{Supp}(M))$ and $\mathfrak p$
is a minimal prime in the support of $M$.
In particular, $M$ has no embedded associated primes.
\end{lemma}

\begin{proof}
By Lemma \ref{lemma-depth-dim-associated-primes} we have
$\text{depth}(M) \leq \dim(R/\mathfrak p)$.
Of course $\dim(R/\mathfrak p) \leq \dim(\text{Supp}(M))$
as $\mathfrak p \in \text{Supp}(M)$ (Lemma \ref{lemma-ass-support}).
Thus we have equality in both inequalities as $M$ is Cohen-Macaulay.
Then $\mathfrak p$ must be minimal in $\text{Supp}(M)$ otherwise
we would have $\dim(R/\mathfrak p) < \dim(\text{Supp}(M))$.
Finally, minimal primes in the support of $M$ are equal to
the minimal elements of $\text{Ass}(M)$
(Proposition \ref{proposition-minimal-primes-associated-primes})
hence $M$ has no embedded associated primes
(Definition \ref{definition-embedded-primes}).
\end{proof}

\begin{definition}
\label{definition-maximal-CM}
Let $R$ be a Noetherian local ring.
A finite module $M$ over $R$ is called a {\it maximal Cohen-Macaulay}
module if $\text{depth}(M) = \dim(R)$.
\end{definition}

\noindent
In other words, a maximal Cohen-Macaulay module over a Noetherian local
ring is a finite module with the largest possible depth over that ring.
Equivalently, a maximal Cohen-Macaulay module over a Noetherian local
ring $R$ is a Cohen-Macaulay module of dimension equal to the dimension
of the ring. In particular, if $M$ is a Cohen-Macaulay $R$-module with
$\Spec(R) = \text{Supp}(M)$, then $M$ is maximal Cohen-Macaulay.
Thus the following two lemmas are on maximal Cohen-Macaulay modules.

\begin{lemma}
\label{lemma-maximal-chain-maximal-CM}
\begin{slogan}
In a local Cohen-Macaulay ring, any maximal chain of prime ideals has
length equal to the dimension.
\end{slogan}
Let $R$ be a Noetherian local ring. Assume there exists a
Cohen-Macaulay module $M$ with $\Spec(R) = \text{Supp}(M)$.
Then any maximal chain of prime ideals $\mathfrak p_0 \subset
\mathfrak p_1 \subset \ldots \subset \mathfrak p_n$
has length $n = \dim(R)$.
\end{lemma}

\begin{proof}
We will prove this by induction on $\dim(R)$. If $\dim(R) = 0$,
then the statement is clear. Assume $\dim(R) > 0$. Then $n > 0$.
Choose an element $x \in \mathfrak p_1$, with $x$ not in
any of the minimal primes of $R$, and in particular
$x \not \in \mathfrak p_0$. (See Lemma \ref{lemma-silly}.)
Then $\dim(R/xR) = \dim(R) - 1$ by Lemma \ref{lemma-one-equation}.
The module $M/xM$ is Cohen-Macaulay over $R/xR$ by
Proposition \ref{proposition-CM-module} and
Lemma \ref{lemma-CM-over-quotient}.
The support of $M/xM$ is $\Spec(R/xR)$ by
Lemma \ref{lemma-support-quotient}.
After replacing $x$ by $x^n$ for some $n$,
we may assume that $\mathfrak p_1$ is an associated prime of $M/xM$, see
Lemma \ref{lemma-inherit-minimal-primes}.
By Lemma \ref{lemma-CM-ass-minimal-support}
we conclude that $\mathfrak p_1/(x)$ is a minimal prime of $R/xR$.
It follows that  the chain
$\mathfrak p_1/(x) \subset \ldots \subset \mathfrak p_n/(x)$
is a maximal chain of primes in $R/xR$.
By induction we find that this chain
has length $\dim(R/xR) = \dim(R) - 1$ as desired.
\end{proof}

\begin{lemma}
\label{lemma-dim-formula-maximal-CM}
Suppose $R$ is a Noetherian local ring. Assume there exists a
Cohen-Macaulay module $M$ with $\Spec(R) = \text{Supp}(M)$. Then for
a prime $\mathfrak p \subset R$ we have
$$
\dim(R) = \dim(R_{\mathfrak p}) + \dim(R/\mathfrak p).
$$
\end{lemma}

\begin{proof}
Follows immediately from Lemma \ref{lemma-maximal-chain-maximal-CM}.
\end{proof}

\begin{lemma}
\label{lemma-localize-CM-module}
Suppose $R$ is a Noetherian local ring. Let $M$ be a Cohen-Macaulay
module over $R$. For any prime $\mathfrak p \subset R$ the
module $M_{\mathfrak p}$ is Cohen-Macaulay over $R_\mathfrak p$.
\end{lemma}

\begin{proof}
We may and do assume $\mathfrak p \not = \mathfrak m$ and $M$ not zero.
Choose a maximal chain of primes $\mathfrak p = \mathfrak p_c \subset
\mathfrak p_{c - 1} \subset \ldots \subset \mathfrak p_1 \subset \mathfrak m$.
If we prove the result for $M_{\mathfrak p_1}$ over $R_{\mathfrak p_1}$,
then the lemma will follow by induction on $c$. Thus we may assume that
there is no prime strictly between $\mathfrak p$ and $\mathfrak m$.
Note that
$\dim(\text{Supp}(M_\mathfrak p)) \leq \dim(\text{Supp}(M)) - 1$
because any chain of primes in the support of $M_\mathfrak p$
can be extended by one more prime (namely $\mathfrak m$) in the
support of $M$. On the other hand, we have
$\text{depth}(M_\mathfrak p) \geq \text{depth}(M) - \dim(R/\mathfrak p) =
\text{depth}(M) - 1$ by Lemma \ref{lemma-depth-localization} and our
choice of $\mathfrak p$.
Thus $\text{depth}(M_\mathfrak p) \geq \dim(\text{Supp}(M_\mathfrak p))$
as desired (the other inequality is Lemma \ref{lemma-bound-depth}).
\end{proof}

\begin{definition}
\label{definition-module-CM}
Let $R$ be a Noetherian ring. Let $M$ be a finite $R$-module.
We say $M$ is {\it Cohen-Macaulay} if $M_\mathfrak p$ is a Cohen-Macaulay
module over $R_\mathfrak p$ for all primes $\mathfrak p$ of $R$.
\end{definition}

\noindent
By Lemma \ref{lemma-localize-CM-module} it suffices to check this
in the maximal ideals of $R$.

\begin{lemma}
\label{lemma-maximal-CM-polynomial-algebra}
Let $R$ be a Noetherian ring. Let $M$ be a Cohen-Macaulay module
over $R$. Then $M \otimes_R R[x_1, \ldots, x_n]$ is a
Cohen-Macaulay module over $R[x_1, \ldots, x_n]$.
\end{lemma}

\begin{proof}
By induction on the number of variables it suffices to prove this for
$M[x] = M \otimes_R R[x]$ over $R[x]$. Let $\mathfrak m \subset R[x]$
be a maximal ideal, and let $\mathfrak p = R \cap \mathfrak m$.
Let $f_1, \ldots, f_d$ be a $M_\mathfrak p$-regular sequence in the maximal
ideal of $R_{\mathfrak p}$ of length $d = \dim(\text{Supp}(M_{\mathfrak p}))$.
Note that since $R[x]$ is flat over $R$ the localization
$R[x]_{\mathfrak m}$ is flat over $R_{\mathfrak p}$.
Hence, by Lemma \ref{lemma-flat-increases-depth}, the sequence
$f_1, \ldots, f_d$ is a $M[x]_{\mathfrak m}$-regular sequence of length $d$
in $R[x]_{\mathfrak m}$. The quotient
$$
Q = M[x]_{\mathfrak m}/(f_1, \ldots, f_d)M[x]_{\mathfrak m} =
M_{\mathfrak p}/(f_1, \ldots, f_d)M_{\mathfrak p}
\otimes_{R_\mathfrak p} R[x]_{\mathfrak m}
$$
has support equal to the primes lying over $\mathfrak p$
because $R_\mathfrak p \to R[x]_\mathfrak m$ is flat and
the support of $M_{\mathfrak p}/(f_1, \ldots, f_d)M_{\mathfrak p}$
is equal to $\{\mathfrak p\}$ (details omitted; hint: follows from
Lemmas \ref{lemma-annihilator-flat-base-change} and
\ref{lemma-support-closed}). Hence the dimension is $1$.
To finish the proof it suffices to find an $f \in \mathfrak m$
which is a nonzerodivisor on $Q$. Since $\mathfrak m$ is
a maximal ideal, the field extension
$\kappa(\mathfrak m)/\kappa(\mathfrak p)$
is finite (Theorem \ref{theorem-nullstellensatz}).
Hence we can find $f \in \mathfrak m$ which
viewed as a polynomial in $x$ has leading
coefficient not in $\mathfrak p$. Such an $f$ acts as
a nonzerodivisor on
$$
M_{\mathfrak p}/(f_1, \ldots, f_d)M_{\mathfrak p} \otimes_R R[x] =
\bigoplus\nolimits_{n \geq 0}
M_{\mathfrak p}/(f_1, \ldots, f_d)M_{\mathfrak p} \cdot x^n
$$
and hence acts as a nonzerodivisor on $Q$.
\end{proof}















